\documentclass[conference]{IEEEtran}
\usepackage{cite}
\usepackage{amsmath,amssymb,amsfonts}
\usepackage{algorithmic}
\usepackage{graphicx}
\usepackage{textcomp}
\usepackage{xcolor}
\usepackage{url}
\usepackage[final]{microtype}
\usepackage{tikz}
\usetikzlibrary{arrows.meta, backgrounds, calc, fit, positioning,
                shadows.blur, shapes.geometric, shapes.misc}

\usepackage{tabularx}
\usepackage{float}
\usepackage{flushend}

\usepackage{calc}

\def\BibTeX{{\rm B\kern-.05em{\sc i\kern-.025em b}\kern-.08em
    T\kern-.1667em\lower.7ex\hbox{E}\kern-.125emX}}
\begin{document}

\title{An Open-source LLM-Assisted System for IT/SRE Triage }

\makeatletter
\newcommand{\linebreakand}{%
  \end{@IEEEauthorhalign}
  \hfill\mbox{}\par
  \mbox{}\hfill\begin{@IEEEauthorhalign}
}
\makeatother



\author{
\IEEEauthorblockN{Kalhar Pandya}
\IEEEauthorblockA{\textit{Khoury College of Computer Sciences} \\
\textit{Northeastern University}\\
Vancouver, Canada \\
pandya.kal@northeastern.edu}
\and
\IEEEauthorblockN{Dawood Sajjadi}
\IEEEauthorblockA{\textit{Fortinet} \\
Consultant}
\and
\IEEEauthorblockN{Maryam Tanha}
\IEEEauthorblockA{\textit{Khoury College of Computer Sciences} \\
\textit{Northeastern University}\\
Vancouver, Canada \\
m.tanha@northeastern.edu}
}


\maketitle
\begin{abstract}
During a production incident, Site Reliability Engineering teams
must SSH into multiple hosts, correlate logs, and reconcile metrics
across monitoring stacks under time pressure. Large Language Models
can accelerate parts of this workflow, yet prior work shows they
struggle with root cause analysis even when given topology and
telemetry data; a single model lacks the structural scaffolding to
reason across it. We present an open-source, hierarchical
multi-agent framework built on LangGraph that provides this
scaffolding through structured decomposition. A Parent Agent drives
iterative investigation cycles, delegating subtasks to parallel
Specialist Agents (ReAct-based workers that run read-only commands
over Docker~exec or SSH) while a Synthesis Agent consolidates
cross-domain findings between cycles. Operators declare their
service configuration in YAML files validated by Pydantic at
startup; adding a new specialist requires only a YAML config, no
Python. A two-layer security model (command allowlist plus output
redaction) prevents agents from modifying infrastructure or leaking
credentials. We evaluate on a five-service mail application
deployed on an Amazon Web Services (AWS) Elastic Compute Cloud (EC2)
instance, injecting six faults across six infrastructure domains.
The agent correctly identifies the root cause in all six cases,
averaging 89~seconds and \$0.27 per investigation on GPT-4.1.
\end{abstract}

\begin{IEEEkeywords}
multi-agent systems, large language models, root cause analysis,
AIOps, site reliability engineering, LangGraph
\end{IEEEkeywords}
\section{Introduction}
When a cloud-native microservice stack fails, the response is
inherently manual. Site Reliability Engineering (SRE) and DevOps
teams open Secure Shell (\texttt{SSH}) sessions to multiple hosts,
correlate log fragments across services that may share nothing but a
request ID, and attempt to isolate root causes from downstream
symptoms under time pressure. Surveys and industrial experience
confirm that this process remains laborious and
error-prone~\cite{zhang2025aiops_survey, chen2024rcacopilot}.
Large Language Models (LLMs) are increasingly applied to parts of
this process, though not as replacements for existing toolchains.
A recent AIOps survey~\cite{zhang2025aiops_survey} recommends
integrating LLMs with established monitoring stacks rather than
building from scratch, and a chaos-engineering
evaluation~\cite{szandala2025aiops_reliability} finds that models
function best as co-pilots alongside human SREs. One finding in particular
shapes our architecture: even when provided with system topology and
oracle telemetry (i.e., ground-truth telemetry with no noise or
missing data), the best-performing LLM solved only 5.37\% of
root cause analysis tasks in the OpenRCA benchmark~\cite{xu2025openrca},
Without additional structural support, models cannot
reason across complex service dependency graphs.
The core limitation is asking a single model to process the full
topology at once.

We address this through multi-agent decomposition over a declared
service topology: rather than asking one model to
reason over the entire stack, we break the problem into
domain-specific subtasks, each grounded in a narrow slice of the
topology, iterate across investigation cycles, and
synthesize findings across agents. The framework grounds each
specialist in a YAML service topology and enforces a read-only
execution model throughout the
investigation~\cite{kholkar2025safeops}. Before each specialist begins, the framework pre-fetches container logs and runs
configured context commands, injecting structured evidence into the agent's opening
message rather than forwarding raw streams that would exhaust a
context window. Structured log representations improve LLM parsing
accuracy~\cite{xu2024divlog}, motivating this pre-processing step.

Our contributions are as follows:
\begin{itemize}
  \item We design and implement an open-source, hierarchical
        multi-agent framework for automated root cause analysis,
        using a Parent--Specialist--Synthesis architecture on
        LangGraph that decomposes investigations into parallel,
        domain-specific subtasks.
  \item We introduce a YAML-based profile system that allows
        operators to describe deployment targets and add specialist
        agents without writing Python code.
  \item We implement a two-layer security model combining a command
        allowlist with output redaction to enforce read-only agent
        execution and prevent credential exposure.
  \item We evaluate the framework on a five-service mail application
        with six injected faults across six infrastructure domains,
        achieving correct diagnosis in all cases at an average of
        89~seconds and \$0.27 per run.
\end{itemize}
\begin{sloppypar}
The source code is available at \url{https://github.com/maryam-tanha/LLM-assisted-Triage}.
\end{sloppypar}

\section{Related Work}

\subsection{Multi-Agent Architectures for IT Operations}
Splitting a complex task across specialized agents outperforms asking a single
model to handle everything. AgentOrchestra~\cite{zhang2025agentorchestra} formalizes this
with the TEA (Tool-Environment-Agent) protocol, where a central planner delegates to
domain-specific sub-agents; the architecture achieves 89.04\% average across difficulty levels on the General AI Assistants (GAIA) benchmark, well above
monolithic baselines. MetaGPT~\cite{hong2024metagpt} assigns each
agent a Standardized Operating Procedure (SOP) modeled on real-world
software team workflows, which constrains output format and reduces
errors during collaboration. In both cases, enforcing role boundaries proved as important as
model capability.

Context management remains a recurring difficulty in multi-agent IT systems.
An AI-Oriented Operations (AOI) system is proposed in~\cite{bai2025aoi}. It
compresses intermediate findings before passing them to a high-level observer, retaining 92.8\%
of critical information at 72.4\% compression, while cutting mean time to resolution by 34.4\%
over the best prior baseline.

Three additional systems inform our design. Triangle~\cite{yu2025triangle} demonstrated over 20\% triage
accuracy improvement (reaching 91.7\% accuracy within five
reassignment hops) in a production deployment serving tens of
millions of users. MasRouter~\cite{yue2025masrouter} routes queries to
appropriately sized models and collaboration configurations when
complexity permits, reducing HumanEval overhead by up to 52\%. SageCopilot~\cite{liao2025sagecopilot}
integrates reflective validation and human-in-the-loop feedback;
its iterative post-generation correction yielded measurable quality
gains.

\subsection{LLM-Based Root Cause Analysis}
Feeding raw diagnostic data directly to an LLM yields low accuracy in production settings. RCACopilot at
Microsoft~\cite{chen2024rcacopilot} addressed this by summarizing structured diagnostic
information before passing it to the prediction model, gaining 0.077 absolute Micro-F1 points
over the unsummarized baseline. FLASH~\cite{zhang2024flash} extends this by decomposing the
full diagnostic lifecycle into four supervised states and learning from historical failure
runs; across 250 real incidents it reached 73.9\% accuracy, and in a separate human evaluation
brought average diagnosis time down to 5.3 minutes against a 90-minute mitigation baseline.

Separating agents by function also affects accuracy.
MA-RCA~\cite{fu2026marca} splits retrieval and validation into distinct agents and reports
95.8\% accuracy on the Nezha cloud-native dataset, a result the authors attribute directly
to eliminating context-switching between hypothesis generation and evidence checking.
D-Bot~\cite{zhou2023dbot} applies tree-search reasoning using Upper Confidence bound applied to Trees (UCT) scoring to database
anomalies, forcing the planner to explore multiple diagnostic paths rather
than committing to a single chain of actions.

Not every scenario suits the same approach. Hy-LIFT~\cite{salman2025hylift} shows that for
5G/6G networks, rule-based labeling handles known fault classes well; the LLM layer is reserved
for novel or ambiguous cases, reaching a macro-F1 of 0.89. At the other end,
an Artificial Intelligence for IT Operations (AIOps) reliability study~\cite{szandala2025aiops_reliability} finds that zero-shot LLMs
misclassify benign traffic spikes as security incidents.

\subsection{Specialist Agent Design and Data Ingestion}
Ashraf and Talavera~\cite{ashraf2025autonomous} observe that
conflicting decisions arise when agent coordination is not
effectively managed, and that assigning specialized roles leads to
more comprehensive evaluation.

Once agents have clear roles, the next question is what data each one should receive.
Dumping raw infrastructure logs into an LLM context reliably fails: token
budgets overflow and important lines get truncated before the model
ever sees them. DivLog~\cite{xu2024divlog} addresses template extraction specifically by sampling a
maximally diverse set of log candidates offline and using them as in-context
examples at inference time. Across 16 LogPAI benchmarks it achieves 98.1\% parsing accuracy,
which makes structured templates a reliable input rather than an aspiration.
Together, these two concerns (clear role boundaries that prevent
cross-domain interference, and structured data ingestion that keeps
each agent's context window focused) form the foundation of our
specialist design.

\subsection{Agent Resilience and Security}
Huang
et al.~\cite{huang2024resilience} inject faulty agents into three topology types and measure
the damage: hierarchical systems lose only 5.5\% accuracy; flat topologies, lacking clear
authority, reach decision paralysis at 10.5\%; linear chains, where each hop re-transmits
errors, degrade by 23.7\%.

On the security side, Kholkar and Ahuja~\cite{kholkar2025safeops} demonstrate a
``policy-as-prompt'' pattern that encodes governance rules into runtime classifiers; their
least-privilege framing directly informs our read-only execution policy.


\section{Methodology}


% ── Colour palette (used only in this figure) ───────────────────
\definecolor{inputbg} {RGB}{232,240,254}
\definecolor{inputbd} {RGB}{160,190,240}
\definecolor{procbg}  {RGB}{220,238,255}
\definecolor{procbd}  {RGB}{130,180,220}
\definecolor{outbg}   {RGB}{230,249,238}
\definecolor{outbd}   {RGB}{100,190,140}
\definecolor{boxbg}   {RGB}{250,252,255}
\definecolor{boxbd}   {RGB}{150,185,225}
\definecolor{parentbg}{RGB}{235,244,255}
\definecolor{parentbd}{RGB}{80,130,200}
\definecolor{synthbg} {RGB}{248,240,255}
\definecolor{synthbd} {RGB}{148,100,210}
\definecolor{agentbd} {RGB}{110,165,220}
\definecolor{arrowcol}{RGB}{70,120,190}
\definecolor{loopcol} {RGB}{220,140,20}
\definecolor{passgreen}{RGB}{60,160,100}
\definecolor{outarrow} {RGB}{50,150,90}

\begin{figure*}[t]
    \centering
    \makebox[\textwidth][c]{%
    \begin{tikzpicture}[
      font=\small\sffamily,
      >=Stealth,
      layer/.style  = {draw, rounded corners=6pt, inner sep=12pt,
                       blur shadow={shadow blur steps=5, shadow xshift=0.5pt,
                                    shadow yshift=-1pt, shadow opacity=15}},
      box/.style    = {draw=boxbd, fill=boxbg, rounded corners=4pt,
                       inner sep=5pt, minimum height=22pt, align=center,
                       blur shadow={shadow blur steps=3, shadow opacity=10}},
      agent/.style  = {draw=agentbd, fill=boxbg, rounded corners=4pt,
                       inner sep=4pt, minimum width=70pt, minimum height=38pt,
                       align=center, text width=66pt,
                       blur shadow={shadow blur steps=3, shadow opacity=10}},
      parent/.style = {draw=parentbd, fill=parentbg, rounded corners=6pt,
                       inner sep=7pt, minimum width=180pt, minimum height=36pt,
                       align=center, line width=1.3pt,
                       blur shadow={shadow blur steps=4, shadow opacity=15}},
      synth/.style  = {draw=synthbd, fill=synthbg, rounded corners=6pt,
                       inner sep=7pt, minimum width=380pt, minimum height=32pt,
                       align=center, line width=1.3pt,
                       blur shadow={shadow blur steps=4, shadow opacity=15}},
      reportbox/.style= {draw=outbd, fill=outbg, rounded corners=6pt,
                       inner sep=7pt, minimum width=220pt, minimum height=38pt,
                       align=center, line width=1.3pt,
                       blur shadow={shadow blur steps=4, shadow opacity=15}},
      arr/.style    = {->, thick, color=arrowcol},
      arrgr/.style  = {->, thick, color=outarrow},
      looparr/.style= {->, very thick, color=loopcol, dashed},
      sublabel/.style = {font=\fontsize{7}{8}\selectfont\sffamily\itshape,
                         color=gray!70},
    ]

    % ── INPUT LAYER ────────────────────────────────────────────────
    \node[box] (src) {
      \textbf{Integration Source}\\[1pt]
      {\fontsize{7}{8}\selectfont PagerDuty / Jira / Manual}
    };

    \node[box, right=12pt of src] (ticket) {
      \textbf{Incident Ticket}\\[2pt]
      {\fontsize{7}{8}\selectfont\textit{title}: [Incident Title]}\\
      {\fontsize{7}{8}\selectfont\textit{severity}: [P1/P2/P3]}\\
      {\fontsize{7}{8}\selectfont\textit{services}: [List]}
    };

    \node[box, right=18pt of ticket] (yaml) {
      \textbf{Profile Directory (YAML)}\\[2pt]
      {\fontsize{7}{8}\selectfont\texttt{profile.yaml} -- services, access}\\
      {\fontsize{7}{8}\selectfont\texttt{agents/*.yaml} -- specialist configs}\\
      {\fontsize{7}{8}\selectfont\texttt{parent.yaml}, \texttt{synthesis.yaml}}
    };

    \begin{pgfonlayer}{background}
      \node[layer, draw=inputbd, fill=inputbg,
            fit=(src)(ticket)(yaml),
            label={[anchor=south west, inner sep=2pt, font=\fontsize{8}{9}\selectfont\sffamily\bfseries,
                    text=gray!60] north west:INPUT LAYER}] (inputlayer) {};
    \end{pgfonlayer}

    \draw[arr] (src.east) -- (ticket.west);

    % ── PROCESSING LAYER ───────────────────────────────────────────
    \node[parent, below=28pt of ticket] (parent) {
      \textbf{\large Parent Agent}\\[2pt]
      {\fontsize{8}{9}\selectfont Plans $\cdot$ Delegates $\cdot$ Converges}
    };

    \node[agent, below=22pt of parent, xshift=-130pt] (logagent) {
      \textbf{Log}\\[-1pt]\textbf{Agent}\\[2pt]
      {\fontsize{7}{8}\selectfont Error patterns,}\\
      {\fontsize{7}{8}\selectfont app logs, \texttt{journalctl}}
    };

    \node[agent, right=8pt of logagent] (rtagent) {
      \textbf{Runtime}\\[-1pt]\textbf{Status Agent}\\[2pt]
      {\fontsize{7}{8}\selectfont CPU, mem, disk,}\\
      {\fontsize{7}{8}\selectfont process state}
    };

    \node[agent, right=8pt of rtagent] (net) {
      \textbf{Network}\\[-1pt]\textbf{Agent}\\[2pt]
      {\fontsize{7}{8}\selectfont Connectivity,}\\
      {\fontsize{7}{8}\selectfont latency, DNS}
    };

    \node[agent, right=8pt of net] (db) {
      \textbf{Database}\\[-1pt]\textbf{Agent}\\[2pt]
      {\fontsize{7}{8}\selectfont Query perf,}\\
      {\fontsize{7}{8}\selectfont locks, replication}
    };

    \node[agent, right=8pt of db, dashed, draw=gray!60, text=gray!80, fill=gray!5] (ext) {
      {\fontsize{8}{9}\selectfont\textit{Custom}}\\[-1pt]
      {\fontsize{8}{9}\selectfont\textit{Specialist}}\\[2pt]
      {\fontsize{7}{8}\selectfont\textit{YAML config or}}\\
      {\fontsize{7}{8}\selectfont\textit{\texttt{register()}}}
    };

    \node[synth, below=20pt of rtagent, xshift=55pt] (synth) {
      \textbf{\normalsize Synthesis Agent}\\[2pt]
      {\fontsize{8}{9}\selectfont
       Cross-Domain Consolidation $\cdot$
       Causal Chain Analysis $\cdot$
       Large Context Window}
    };

    \begin{pgfonlayer}{background}
      \node[layer, draw=procbd, fill=procbg,
            fit=(parent)(logagent)(rtagent)(net)(db)(ext)(synth),
            label={[anchor=south west, inner sep=2pt, font=\fontsize{8}{9}\selectfont\sffamily\bfseries,
                    text=gray!60] north west:PROCESSING LAYER}] (proclayer) {};
    \end{pgfonlayer}

    % ── OUTPUT LAYER ───────────────────────────────────────────────
    \node[reportbox, below=36pt of synth] (report) {
      \textbf{\normalsize Final RCA Report}\\[2pt]
      {\fontsize{8}{9}\selectfont
       Root cause $\cdot$ Evidence $\cdot$
       Contributing factors $\cdot$ Recommended actions}
    };

    \begin{pgfonlayer}{background}
      \node[layer, draw=outbd, fill=outbg,
            fit=(report),
            inner sep=16pt,
            label={[anchor=south west, inner sep=2pt, font=\fontsize{8}{9}\selectfont\sffamily\bfseries,
                    text=gray!60] north west:OUTPUT LAYER}] (outlayer) {};
    \end{pgfonlayer}

    % ── ARROWS ─────────────────────────────────────────────────────

    % Input → Parent
    \draw[arr] (ticket.south)
        -- node[left, sublabel]{incident + config}
        (parent.north);

    % YAML → Parent
    \coordinate (yamldown) at ($(yaml.south) + (0,-10pt)$);
    \draw[arr, dashed, color=gray!60]
        (yaml.south)
        -- (yamldown)
        -- node[above, sublabel, xshift=-28pt]{Pydantic-validated}
        (parent.east |- yamldown)
        -- (parent.east);

    % Parent → Specialists (fan-out)
    \coordinate (fanout) at ($(parent.south) + (0,-12pt)$);
    \draw[arr] (parent.south) -- node[left, sublabel, yshift=-1pt]{\texttt{Send()} fan-out} (fanout) -| (logagent.north);
    \draw[arr] (fanout) -| (rtagent.north);
    \draw[arr] (fanout) -| (net.north);
    \draw[arr] (fanout) -| (db.north);

    % Specialists → Synthesis (fan-in)
    \coordinate (fanin) at ($(synth.north) + (0,12pt)$);
    \draw[arr] (logagent.south) -- ++(0,-6pt) -| (fanin);
    \draw[arr] (rtagent.south)  -- ++(0,-6pt) -| (fanin);
    \draw[arr] (net.south)      -- ++(0,-6pt) -| (fanin);
    \draw[arr] (db.south)       -- ++(0,-6pt) -| (fanin);
    \draw[arr] (fanin) -- node[right, sublabel]{\texttt{SpecialistFinding}} (synth.north);

    % Synthesis → Parent (investigation loop, left side)
    \coordinate (loopX) at ($(logagent.west) + (-20pt, 0)$);
    \draw[looparr] (synth.west) -- (synth.west -| loopX)
        -- node[left, sublabel, align=center, fill=procbg, inner sep=2pt]{%
             \textbf{Investigation}\\
             \textbf{Loop}\\
             (max $N$ cycles)}
        (parent.west -| loopX) -- (parent.west);

    % Parent → Final Report (when concluding)
    \coordinate (concludeRight) at ($(proclayer.east) + (18pt, 0)$);
    \coordinate (concludeMidY) at ($(report.north) + (0, 25pt)$);
    \draw[arrgr, very thick, rounded corners=6pt]
        (parent.east)
        -- node[above, sublabel, pos=0.8]{\texttt{write\_rca\_conclusion}}
        (parent.east -| concludeRight)
        -- (concludeRight |- concludeMidY)
        -- (report.north |- concludeMidY)
        -- (report.north);

    % ── Title ──────────────────────────────────────────────────────
    \node[above=6pt of inputlayer,
          font=\normalsize\sffamily\bfseries]
      {System Overview: Hierarchical Multi-Agent RCA Architecture};

    \end{tikzpicture}%
    }% end makebox
    \caption{Multi-agent RCA architecture. The Parent Agent fans out
      investigation subtasks in parallel via \texttt{Send()} and iterates
      until the accumulated evidence supports a conclusion or the cycle
      limit is reached. The dashed orange path shows the investigation
      loop; the green path marks the conclusion route to the final report.}
    \label{fig:architecture}
    \end{figure*}


Fig.~\ref{fig:architecture} shows the full system. We built it as a
three-tier LangGraph graph: a Parent Agent that plans and decides, a
pool of Specialist Agents that gather evidence in parallel, and a
Synthesis Agent that correlates findings across domains before the
next cycle. Specialist workers report exclusively through the orchestrator because
flat and linear topologies suffer higher accuracy loss under
agent failure~\cite{huang2024resilience}.
Each product is described in its own profile directory of YAML
files, validated by Pydantic at load time; pointing the framework at
a different environment requires editing configuration, not code.

\subsection{Service Configuration}

The goal is to let operators describe their infrastructure once, in a
format they already understand, without writing any agent-specific code. The
result is a \emph{profile directory} per product, validated by Pydantic
at startup:

\begin{itemize}
  \item \texttt{profile.yaml} -- product metadata, access method
        (\texttt{docker\_exec} or \texttt{ssh}), and a
        \texttt{services} block where each entry captures what a
        healthy instance looks like (\texttt{expected\_behavior}),
        what patterns signal trouble (\texttt{known\_failures}),
        and which shell commands collect an initial evidence snapshot
        (\texttt{context\_commands}).
  \item \texttt{agents/*.yaml} -- one file per specialist, each
        declaring an \texttt{agent\_type}, a domain-specific
        \texttt{system\_prompt}, routing guidance
        (\texttt{when\_to\_use}, \texttt{do\_not\_use}),
        optional agent-level \texttt{context\_commands}, and a
        \texttt{gather\_docker\_host\_context} flag that controls
        whether Docker-level metadata is collected before the
        specialist begins.
  \item \texttt{parent.yaml} and \texttt{synthesis.yaml} -- customizable
        system prompts for the Parent and Synthesis agents, allowing
        operators to adjust orchestration behavior per product
        without modifying framework code.
\end{itemize}

At runtime the Parent Agent reads the specialist routing hints
verbatim when choosing whom to assign. Nothing is hardcoded in the
planner prompt; the entire profile is loaded at invocation time, so
onboarding a new product amounts to filling in a directory.

Specialists can also be introduced without writing Python at all. A
YAML agent configuration file is enough: a generic
\texttt{YAMLSpecialist} bridge class picks it up at startup,
builds the corresponding ReAct agent, and wires the node into the
graph. An operator familiar with Kafka internals, for instance, can
contribute a specialist by writing a system prompt and a few routing
hints; the framework handles registration and graph wiring.

\subsection{Parent Agent and Planning}

The Parent Agent is the only node in the graph that can change the
investigation's direction. At each cycle it receives the incident
summary, the full YAML topology, and the accumulated
\texttt{cumulative\_history} from all prior cycles. We bind the LLM
with \texttt{tool\_choice=``any''} to ensure it always produces a
typed action rather than free text. This replanning pattern follows multi-agent RCA
designs that explicitly separate retrieval from synthesis~\cite{fu2026marca, zhou2023dbot}.

The LLM has exactly two choices. \texttt{create\_subtasks} emits a
list of \texttt{Subtask} objects, each specifying a target service,
container, natural-language description, a falsifiable hypothesis, and
the agent to assign. The LangGraph router converts these into
\texttt{Send()} fan-out dispatches that run in parallel.
\texttt{write\_rca\_conclusion} closes the cycle, producing a
structured report with root cause, contributing factors, evidence
summary, and recommended actions.

We chose not to rely on the agent voluntarily concluding: once the zero-indexed \texttt{current\_cycle} counter reaches
\texttt{MAX\_CYCLES} (default:~3, yielding up to four investigation
rounds), the framework hard-wires \texttt{write\_rca\_conclusion}
as the only permitted tool call; the loop cannot run indefinitely.

Not every incident requires autonomous planning from scratch. The
framework therefore supports a \emph{manual investigation mode}: an
operator who already suspects a specific service can pre-populate the
subtask list before launching the graph. If subtasks exist at
cycle~0, the Parent Agent skips the LLM call and fans them out
directly. Any subsequent cycles fall back to normal LLM-driven
planning, so the autonomous loop remains available even when the
first round is hand-directed.

\subsection{Specialist Agents}

A specialist is a LangChain Reasoning + Acting (ReAct) agent with a deliberately narrow
scope. Its only tool is \texttt{run\_command}; its system prompt
encodes domain knowledge for one area (logs, runtime resources,
network connectivity, container configuration, etc.) and nothing else.
Empirically, narrowing the system prompt reduced
cross-domain confusions: a log specialist with Redis knowledge
in its prompt issued Redis commands during unrelated log
investigations.

Before the ReAct loop begins, the framework pre-runs the service's
\texttt{context\_commands} and, in Docker mode, tails the container
logs. Both are injected into the agent's opening message, giving it
a concrete evidence baseline without spending any tool-call budget; log output is
pre-processed and structured before injection, following evidence
that structured templates improve LLM parsing~\cite{xu2024divlog}.
When the per-agent flag \texttt{gather\_docker\_host\_context} is
set, the framework additionally collects \texttt{docker inspect}
(full container configuration and state), \texttt{docker stats
-{}-no-stream} (a point-in-time Central Processing Unit (CPU), memory, network, and Input/Output (I/O)
snapshot), and \texttt{docker events} filtered to the preceding
24~hours (daemon-level signals: Out-of-Memory (OOM) kills, restarts, health-check
failures). This host-level collection is necessary because application logs
miss certain failure modes entirely: a container killed by the
kernel OOM reaper, for example, leaves no application-level trace.

Regardless of how evidence is gathered, the loop runs until the agent emits a structured final answer or
hits \texttt{MAX\_ITERATIONS} (default:~10), at which point a
low-confidence \texttt{SpecialistFinding} is returned so the cycle
can continue rather than crash.

Commands run through one of two backends, both presenting the same
interface to the agent. \textit{Docker exec} calls
\texttt{docker exec <container> sh -c <cmd>} on the local host and
is the default for demo deployments. \textit{SSH} uses a pooled
paramiko client keyed by \texttt{(host, port, username)}, reusing
open transport sessions to avoid repeated handshake overhead in
production environments.

Adding a specialist requires no changes to the orchestration layer.
A new Python module calls \texttt{register()} at import time; the graph
builder wires the new node in at startup and the Parent Agent sees it
in its routing prompt the next invocation. As noted in
Section~III-A, a specialist can also be defined purely through YAML:
the \texttt{YAMLSpecialist} class reads the config, assembles the
prompt and tool bindings, and registers the node, so operators can
extend coverage without writing Python.

\subsection{Synthesis}

After each cycle's specialists finish, the Synthesis Agent reads new
\texttt{SpecialistFinding} objects tracked via a \texttt{findings\_offset} pointer to
avoid reprocessing prior cycles. Each finding conforms to a fixed typed schema, keeping
synthesis predictable. The agent
produces a \texttt{CycleSummary}: a narrative paragraph,
\texttt{KEY\_FINDINGS} bullets, and \texttt{RECOMMENDATIONS}. We give
it no tool access deliberately; its job is cross-domain correlation,
noticing for example that a memory spike on the Redis container
explains the connection-pool exhaustion a log agent reported on the
API service. Keeping it tool-free means reasoning over the full
evidence window without branching into further data collection.
The resulting summary feeds into \texttt{cumulative\_history}, which
the Parent Agent reads at the start of every subsequent cycle.

Once the Parent Agent calls \texttt{write\_rca\_conclusion}, the
structured report (root cause, contributing factors, evidence
summary, and recommended actions) goes directly to the operator.

\subsection{Security Model}

We treat agents as untrusted code paths. Every command, regardless of
which specialist issues it, passes through two sequential checks before
execution and before its output reaches the LLM.

The \textit{command allowlist} applies a deny-first rule: a command is
rejected the moment it matches any blocked pattern (destructive
operations such as \texttt{rm} and \texttt{chmod}; shell redirection;
code injection via \texttt{| bash}; and outbound network tools like
\texttt{curl} and \texttt{nc}). Only commands whose prefix appears on
the explicit allow-list (log readers, system status utilities, and
service Command-Line Interfaces (CLIs) like \texttt{redis-cli} and \texttt{psql}) proceed to
execution. Rejected commands return \texttt{"BLOCKED: <reason>"} so
the agent can adapt rather than hang.

The \textit{output redactor} scrubs every response before it leaves
the execution layer. It strips AWS access keys, PEM private-key
blocks, bearer tokens, credential assignments
(\texttt{api\_key=}, \texttt{password=}), email addresses, full Internet Protocol version 4 (IPv4)
addresses (partially preserved for subnet context), and long base64
strings. The redactor runs on every execution path, including the
host-side \texttt{docker logs} fetch, so no raw secrets ever appear
in the LLM's context.

Both checks run at every node without exception. The least-privilege
principle~\cite{kholkar2025safeops} holds across the entire system:
agents can read diagnostic state but cannot alter it or exfiltrate
credentials.

\subsection{Data Privacy and Provider Governance}

Diagnostic data flowing through a live investigation (log excerpts,
service configurations, error traces, topology details) is operationally
sensitive. Preventing this data from
entering model training pipelines is a hard design constraint.

All LLM inference in our framework is routed through the OpenRouter
API rather than through direct provider endpoints. OpenRouter is a
routing proxy that dispatches each request to a backend model provider
but, by default, retains only request metadata (token counts, latency,
model identifier) and not the prompt or completion
content~\cite{openrouter2025privacy}. Two opt-in flags exist for
users who want logging or data sharing, but both are disabled in our
deployment.

The more consequential setting is \texttt{provider.data\_collection},
which can be set to \texttt{"deny"} either per-request or at the
account level. In our deployment we enable this restriction
account-wide, directing OpenRouter to route calls \emph{only} to
providers whose terms of service prohibit collecting or training
on user data. If no
compliant provider is available for the requested model, the call
fails rather than silently falling back to one that might retain
inputs. A stricter option, \texttt{provider.zdr} (zero data
retention), is available for deployments that require providers to hold
nothing beyond the duration of a single request-response cycle.

The providers we tested against, primarily OpenAI (GPT-4.1) via
OpenRouter, independently enforce a no-training default for API
traffic. Anthropic, Mistral, and Google apply similar policies on
their paid API tiers. These guarantees are contractual rather than
technical: routing away from non-compliant providers is the
enforcement mechanism, which is standard practice for any API proxy
architecture. Operators deploying our framework in regulated
environments should verify that their chosen model and provider
combination meets their specific compliance requirements, since
retention windows and audit-log practices vary across vendors.


\section{Experiments and Results}

\subsection{Experimental Setup}

We evaluated the framework on a five-service mail application
deployed on an AWS EC2 instance (Ubuntu) via Docker Compose. The
stack comprises Postfix (Simple Mail Transfer Protocol (SMTP) transport),
Dovecot (Internet Message Access Protocol (IMAP)), Roundcube
with Apache and PHP (webmail front-end), PostgreSQL (address-book
and session storage), and Redis (session cache). This target was chosen over the simpler voting
application used during development because it exercises a wider
range of infrastructure domains (mail transport, web serving,
database connectivity, and Domain Name System (DNS) resolution) within a single
deployment.

Load was generated with Locust, configured for 20~concurrent virtual
users over 90-second windows. Each Locust user performed a mix of
SMTP sends, IMAP logins and mailbox operations, and Roundcube Hypertext Transfer Protocol (HTTP)
workflows (login, inbox listing, compose, send). The framework ran
with default settings: \texttt{MAX\_CYCLES}$=$3,
\texttt{MAX\_ITERATIONS}$=$10 per specialist, and four specialist
agents (Log, Runtime Status, Docker Specs, Network). All inference
used OpenAI's GPT-4.1 model via OpenRouter.

We report four metrics per experiment: whether the agent correctly
identified the injected root cause (binary accuracy), the number of
investigation cycles consumed, wall-clock elapsed time from graph
entry to final report, and the estimated cost in US dollars based on
OpenRouter per-token pricing.

\subsection{Fault Injection}

Each experiment introduces a single-parameter misconfiguration or
corruption into one service while leaving the rest of the stack
healthy. The faults were designed to mimic realistic operator
errors (a connection limit set too low, a memory ceiling that
triggers the OOM killer) rather than synthetic crash
scenarios. Table~\ref{tab:faults} summarizes the six injected
faults, spanning six infrastructure domains: mail protocol,
database, web server, mail transport, DNS, and OS-level memory.

Every fault was injected via \texttt{docker~exec} or a host-side
command, and a Locust run captured the resulting failure signature
before launching the RCA agent. After each experiment the fault was
reversed to restore the baseline.

\begin{table}[htbp]
\centering
\caption{Injected faults and affected services.}
\label{tab:faults}
\renewcommand{\arraystretch}{1.15}
\small
\begin{tabularx}{\columnwidth}{c l X l}
\hline
\textbf{ID} & \textbf{Domain} & \textbf{Fault} & \textbf{Affected} \\
\hline
1 & Mail protocol & Dovecot \path{mail_max_userip_connections=1} & IMAP \\
2 & Database      & PostgreSQL \path{max_connections=4} & Roundcube \\
3 & Web server    & PHP \path{memory_limit=10M} & Roundcube \\
4 & Mail server   & Postfix \path{smtpd_client_message_rate_limit=1} & SMTP \\
5 & DNS           & Corrupt \path{/etc/resolv.conf} in mailserver & Outbound mail \\
6 & OS / Memory   & Container hard memory limit 48\,MB (OOM) & Roundcube \\
\hline
\end{tabularx}
\end{table}

\subsection{Results}

The agent identified the correct root cause in all six experiments.
Table~\ref{tab:results} shows the per-experiment breakdown.

\begin{table}[htbp]
\centering
\caption{RCA agent performance across six fault injections.}
\label{tab:results}
\renewcommand{\arraystretch}{1.15}
\begin{tabularx}{\columnwidth}{c c r r c}
\hline
\textbf{ID} & \textbf{Cycles} & \textbf{Time (s)} & \textbf{Cost (\$)} & \textbf{Correct} \\
\hline
1 & 4 & 136.5 & 0.49 & \checkmark \\
2 & 1 &  30.0 & 0.03 & \checkmark \\
3 & 2 &  74.8 & 0.14 & \checkmark \\
4 & 1 &  44.9 & 0.04 & \checkmark \\
5 & 3 & 155.2 & 0.58 & \checkmark \\
6 & 2 &  95.3 & 0.32 & \checkmark \\
\hline
\multicolumn{2}{l}{\textit{Mean}} & 89.5 & 0.27 & 6/6 \\
\hline
\end{tabularx}
\end{table}

Investigation times ranged from 30~seconds (EXP-02, a single-cycle
diagnosis) to just over two and a half minutes (EXP-05, three
cycles). Token costs stayed between \$0.03 and \$0.58 per run, with
an average of \$0.27.

\subsection{Analysis}

Fault indirection correlates directly with cycle count. Faults that left an unambiguous
signature in the affected service's own logs or configuration
(the PostgreSQL connection pool in EXP-02, the Postfix rate-limit
rejection code in EXP-04) were resolved in a single cycle. The
agent read the error, confirmed the configuration, and concluded.
Faults whose symptoms surfaced one layer away from the actual cause
required more work. In EXP-05, for instance, DNS resolution failures
appeared in Postfix delivery logs, but the root cause was a corrupted
\texttt{/etc/resolv.conf} inside the container. The agent spent its
first cycle investigating Postfix configuration (which was clean),
broadened its network investigation in cycle~2, and only in cycle~3
checked the resolver file directly. This three-cycle progression exercised the iterative narrowing
that the synthesis loop provides.

EXP-06 highlights why host-level container metadata matters. An OOM
kill leaves no trace in application logs; the process simply
disappears. The Docker Specs specialist, running with
\texttt{gather\_docker\_host\_context} enabled, retrieved
\texttt{docker inspect} (which reported \texttt{OOMKilled:~true} and
an incrementing restart count) and \texttt{docker stats} (which
showed memory pinned at 100\% of the 48\,MB limit). Without these
signals the agent would have been limited to observing intermittent
502 errors with no clear cause.

Breadth-first specialist dispatch also contributed to the results. In
every experiment the Parent Agent sent subtasks to all four
specialists in cycle~1, even when the incident description pointed
toward a single service. Specialists that found nothing still
contributed by ruling out competing hypotheses: confirming that
Redis, DNS, or host resources were healthy allowed the Synthesis
Agent to rank the real finding higher in its hypothesis tier.

At an average of \$0.27 per investigation, token cost is
negligible relative to the engineering time a manual triage
typically consumes. The most expensive run (EXP-05, \$0.58)
corresponds to the three-cycle investigation; single-cycle cases
came in under five cents.


\section{Conclusion and Future Directions}

We presented an open-source, hierarchical multi-agent framework for
automated root cause analysis in microservice environments. The
central result is that decomposing an investigation into
domain-specific agents, each grounded in a narrow slice of a
declared service topology, succeeds where prior single-model
approaches have not~\cite{xu2025openrca}. Iterative cycles with
cross-domain synthesis were necessary for faults whose symptoms
surfaced in a different service than the actual cause. The
deny-first security model kept agents useful throughout without
giving them the ability to modify infrastructure or leak
credentials.

The framework's profile system lets operators describe new deployment
targets entirely in YAML (services, specialist agents, and
orchestrator prompts) without writing Python. On our five-service
mail application, six fault injections spanning six infrastructure
domains were all diagnosed correctly, averaging 89~seconds and
\$0.27 per run on GPT-4.1. Direct configuration faults resolved in
one cycle; faults whose symptoms appeared in a different service's
logs took up to three, exercising exactly the iterative narrowing the
architecture was designed for.

Several limitations remain. We tested a single model (GPT-4.1) on
one infrastructure target with controlled fault injections, not
organic incidents. Different architectures or smaller models may
trade accuracy for cost in ways we have not measured. Production
environments are noisier, faults compound, and the ``correct'' root
cause is sometimes a matter of judgment rather than a single
misconfigured parameter. A broader evaluation covering multiple
models, Kubernetes-native deployments, and real incident
data would strengthen the claims considerably.

The most immediate next step is an automated evaluator agent that
scores draft reports on accuracy, completeness, and actionability,
following the review pattern validated by
SageCopilot~\cite{liao2025sagecopilot}. We also plan to adopt the
Model Context Protocol~\cite{aws2025mcp_sre} for dynamic tool
discovery, so specialists can pick up new monitoring backends
without per-integration code. Longer-term, a human-in-the-loop
command approval mode would serve environments where full autonomy
is not yet appropriate: the agent proposes a command, the operator
approves or rejects it, and the framework records the decision for
future policy learning.


\section*{Acknowledgment}
This work was conducted as part of a Research Apprenticeship at
Khoury College of Computer Sciences, Northeastern University, under
the supervision of Dr.\ Maryam Tanha. The authors thank Khoury
College for providing the AWS and OpenRouter credits used in the
experimental evaluation.

\bibliographystyle{ieeetr}
\bibliography{references}

\end{document}
